% 分部积分法（综述）
% license CCBYSA3
% type Wiki

本文根据 CC-BY-SA 协议转载翻译自维基百科\href{https://en.wikipedia.org/wiki/Integration_by_parts}{相关文章}

在微积分中，更一般地说，在数学分析中，分部积分（integration by parts，或称偏积分 partial integration）是一种方法，用来将函数乘积的积分，转化为它们的导数与原函数乘积的积分形式。它常用于把一个函数乘积的不定积分转化为另一个更容易求解的不定积分。这个公式可以被看作是微分中乘积法则的积分版本；事实上，它正是由乘积法则推导而来的。

分部积分公式表述如下：
$$
\begin{aligned}
\int_{a}^{b} u(x)\,v'(x)\,dx &= \Big[u(x)v(x)\Big]_{a}^{b} - \int_{a}^{b} u'(x)\,v(x)\,dx \\
&= u(b)v(b) - u(a)v(a) - \int_{a}^{b} u'(x)\,v(x)\,dx .
\end{aligned}~
$$
或者，设$u = u(x), \quad du = u'(x)\,dx, \quad v = v(x), \quad dv = v'(x)\,dx$,则公式可以写得更为简洁：
$$
\int u\,dv = uv - \int v\,du .~
$$
前一个公式写作定积分形式，而后一个公式则写作不定积分形式。若对后者加上适当的积分限，应当能得到前者，但后者并不必然与前者等价。

数学家布鲁克·泰勒于1715 年首次发表了分部积分的思想。对于 Riemann–Stieltjes 积分与 Lebesgue–Stieltjes 积分，存在更一般化的分部积分形式。而在序列的离散情形中，其对应的类比被称为分部求和。
\subsection{定理}
\subsubsection{两个函数的乘积}
该定理可以如下推导。对于两个连续可微函数$u(x)$和$v(x)$，乘积法则表明：

$$
(u(x)v(x))' = u'(x)v(x) + u(x)v'(x).~
$$

对等式两边关于 $x$ 积分：

$$
\int (u(x)v(x))'\,dx = \int u'(x)v(x)\,dx + \int u(x)v'(x)\,dx,~
$$

注意到不定积分就是原函数，便得到：

$$
u(x)v(x) = \int u'(x)v(x)\,dx + \int u(x)v'(x)\,dx,~
$$

其中省略了积分常数。由此得到分部积分公式：

$$
\int u(x)v'(x)\,dx = u(x)v(x) - \int u'(x)v(x)\,dx,~
$$

或者用微分形式表示：

$$
du = u'(x)\,dx, \quad dv = v'(x)\,dx,~
$$

则有：

$$
\int u(x)\,dv = u(x)v(x) - \int v(x)\,du.~
$$

这应当理解为函数之间的相等式，两边均可加上一个未指定的常数。若在 $x=a$ 与 $x=b$ 两点取差，并应用微积分基本定理，就得到定积分形式：

$$
\int_{a}^{b} u(x)v'(x)\,dx = u(b)v(b) - u(a)v(a) - \int_{a}^{b} u'(x)v(x)\,dx.~
$$
原始积分$\int u v'\,dx$包含导数$v'$；要应用该定理，必须先找到$v$，即$v'$的一个原函数，然后再计算所得积分$\int v u'\,dx$。
